\section{Brief review of the literature}

\subsection*{Theoretical review}

Equity valuation provides a basis for examining reported accounting measures in
relation to share prices. Miller and Modigliani's dividend-irrelevance result
holds under restrictive assumptions, including fixed investment policy and
frictionless capital markets; under those assumptions, changing the payout
does not by itself change firm value \citep{miller1961}. This result is a
benchmark rather than a direct prediction that a higher observed dividend per
share must be associated with a higher market price in Nepal.

Signaling theory offers a conditional reason why accounting disclosures and
dividend decisions may be associated with price. Ross (1977) models financial
structure choices as signals when managers have information unavailable to
investors. Bhattacharya (1979) develops a separate model in which dividends can
signal managers' expectations of future profitability under specified
information and signaling-cost assumptions. These theories do not establish
that every dividend increase signals stronger performance; they motivate
examining DPS alongside the other reported measures \citep{ross1977,
bhattacharya1979}.

Ohlson's valuation model relates market value to current and future earnings,
book values, dividends, and other information, with the clean-surplus relation
providing an accounting foundation \citep{ohlson1995}. In the present study,
EPS represents earnings on a per-share basis, BVPS represents book equity on a
per-share basis, DPS records the reported distribution for the fiscal year,
and ROA represents reported asset profitability. EPS and ROA share an earnings
numerator, while DPS and BVPS are also related to earnings and retained
earnings. The regressors are therefore not assumed to be statistically
independent; their overlap is one reason to assess multicollinearity and to
interpret conditional coefficients cautiously.

Several Nepalese studies include the price-earnings ratio (P/E) alongside
market price per share (MPS) as an explanatory variable \citep{bhattarai2014,
chhetri2023}. Because P/E is calculated as MPS divided by EPS, it is
mathematically coupled to this study's dependent variable. This proposal
therefore excludes P/E from its regression. This is a specification choice for
the present model, not a claim that all prior studies using P/E are invalid.

\subsection*{Empirical review}

International studies provide context but are not direct evidence for the
Nepalese banking sample. Research on NSE 100 firms in India and listed banks in
Jordan examines combinations of earnings, book value, dividends, and valuation
multiples, with findings depending on the market, sample, and specification
\citep{malhotra2013,almumani2014}. Earlier Nepalese bank studies by Bhattarai
(2014) and Pradhan and Dahal (2016) also report associations between accounting
fundamentals and share prices, but their observation periods end before the
early-2023 merger wave \citep{bhattarai2014,pradhan2016}.

More recent Nepalese studies do not yield a uniform result for all variables.
Gyawali (2022), using ten banks over 2017--2021, reports a positive but
statistically insignificant ROA coefficient. Chhetri (2023) studies 13 banks
over 2012--2022 using pooled cross-sectional analysis; the paper reports
positive, statistically significant results for BVPS and ROA and includes P/E
in a model with MPS as the outcome. Joshi et al. (2023) report a positive BVPS
association in their commercial-bank sample. These estimates differ in sample,
period, variable construction, and model specification, so they should not be
treated as estimates of one common effect \citep{gyawali2022,chhetri2023,
joshi2023}.

Studies with periods extending through FY 2023/24 provide evidence from the
post-consolidation period. Pant et al. (2024), using seven purposively selected
banks, report a negative, statistically insignificant BVPS coefficient.
Chaudhari (2025), using ten banks, reports a significantly negative ROA
coefficient. The abstract of Timalsina (2026), which covers 13 banks through
FY 2023/24, reports no statistically significant ROA association. These
studies show that post-consolidation-period evidence exists, but their results
remain mixed across samples, periods, and specifications \citep{pant2024,
chaudhari2025,timalsina2026}.

Merger evidence must also be distinguished by period and outcome. Adhikari,
Kavanagh, and Hampson (2023) compare financial ratios before and after mergers
for seven commercial banks involved in transactions from 2013 to 2020. Their
reported 46.01\% fall in mean MPS is specific to that sample and period; it is
not an estimate for the early-2023 mergers \citep{adhikari2023}. Other work
examines post-merger financial performance or governance rather than the
disclosure-timed share-price outcome used here \citep{adhikari2023cbr,
bhattarai2024,dongol2024}.

\subsection*{Research gap}

The reviewed literature includes studies with observations through FY
2023/24, so this proposal does not claim that post-consolidation evidence is
absent. Instead, it asks what bivariate and conditional associations are
observed in a target of five purposively selected bank series over FY 2076/77--2081/82
when MPS is measured using one explicit rule: the NEPSE closing price on the
first trading day after the corresponding audited annual financial statements
are publicly disclosed through NEPSE. This rule makes the proposed price
observation reproducible; it is not, by itself, evidence that the design is
novel. The study does not test merger effects or resolve conflicting findings
across different studies.

The analysis estimates bivariate associations and a joint pooled-OLS
specification for EPS, DPS, BVPS, and ROA with post-disclosure MPS. It excludes
P/E because that ratio is constructed from MPS and EPS. Findings will be
interpreted for the sampled bank series only, with the small sample and
changes in reporting boundaries around mergers treated as limitations.
